\slajdid{09_2-s022}
% element: 09_2-s022:e01
% element: 09_2-s022:e02
% element: 09_2-s022:e03
% element: 09_2-s022:e04
\begin{frame}[label=09_2-s022]{Pseudo NMOS: naponski nivoi i odnos $\beta$}
\setlength{\abovedisplayskip}{3pt}\setlength{\belowdisplayskip}{3pt}
\begin{columns}[T,onlytextwidth]\column{.36\textwidth}\centering {\RazaviSetup
\begin{circuitikz}[x=1cm,y=1cm,font=\fontsize{10}{12}\selectfont]
\node[rz nmos] (N) at (.28,.65) {};
\node[rz pmos] (P) at (.28,2.2) {};
\coordinate (O) at (N.D |- 0,1.43);
\draw[wire] (N.D)--(O)--(P.D);
\draw[wire] (N.S)--++(0,-.20) coordinate (GND);
\RazaviGround{GND}
\draw[wire] (P.S)--++(0,.22) coordinate (VDD);
\RazaviSupply{VDD}{V_{DD}}
\draw[wire] (O)--++(.60,0) node[rz terminal]{} node[right=3pt]{$V_O$};
\fill (O) circle(1.15pt);
\node[right=3pt] at (N.east) {$\mathrm{N}$};
\node[right=3pt] at (P.east) {$\mathrm{P}$};
\draw[wire] (N.G)--++(-.30,0) node[rz terminal]{} node[left=3pt]{$V_I$};
\draw[wire] (P.G)--++(-.16,0)--++(0,-.18) coordinate (PG);
\RazaviGround{PG}
\end{circuitikz}}
\par\vspace{1mm}\centering \begin{tikzpicture}[x=.64cm,y=.56cm,font=\sitneoznake,>=Latex]
\draw[->](-.15,0)--(4.5,0)node[right]{$V_I$};\draw[->](0,-.15)--(0,3.4)node[above]{$V_O$};
\draw(0,3.05)node[left]{$V_{OH}$}--(.25,3.05)..controls(1.3,3.05)and(1.62,2.68)..(1.85,2.2)..controls(2.2,1.4)and(1.98,.65)..(2.65,.37)..controls(3.0,.28)and(3.4,.22)..(3.7,.2);
\draw[dashed](0,.2)node[left]{$V_{OL}$}--(3.7,.2)--(3.7,0)node[below]{$V_{DD}$};\draw[dashed](1.6,2.58)--(1.6,0)node[below]{$V_{IL}$};\draw[dashed](2.5,.43)--(2.5,0)node[below]{$V_{IH}$};
\draw(1.33,2.95)--(1.98,2.3);\node[right]at(1.6,2.85){$a=-1$};\draw(2.22,.74)--(2.84,.12);\node[right]at(2.5,.82){$a=-1$};
\draw[dashed](0,0)--(2.04,2.04)node[right]{$V_S$};
\end{tikzpicture}
\column{.60\textwidth}$|V_{Tp}|=V_{Tn}$; zbog održivosti je $V_{OL}<V_{Tn}$.\[\frac{k_n}{k_p}=\beta>\frac{V_{DD}-V_{Tn}}{2V_{Tn}}\]
Obavezno $V_{O(IH)}<V_{IL}$; poželjno $V_{O(IH)}<V_{Tn1}$.\[(V_{DD}-V_{Tn})\sqrt{\frac{k_p}{3k_n}}=(V_{DD}-V_{Tn})\sqrt{\frac1{3\beta}}<V_{Tn}\]
\[\beta>\frac13\left(\frac{V_{DD}-V_{Tn}}{V_{Tn}}\right)^2\]
\end{columns}
\end{frame}
